\chapter{Aritmetika kompjuterike dhe analiza e gabimeve}
\begin{qellime}
Studenti përshkruan numrat me pikë të lëvizshme; dallon gabimin absolut e relativ, gabimin e përparmë e të prapmë; analizon kushtëzimin dhe qëndrueshmërinë; njeh anulimin, humbjen e rëndësisë dhe grumbullimin e gabimeve.
\end{qellime}

\section{Përfaqësimi me saktësi të fundme}
Një numër normal binar përfaqësohet si
\begin{equation}
 x=(-1)^s(1.b_1b_2\ldots b_p)_2\,2^e.
\end{equation}
Në standardin IEEE 754, formati double ka 53 bita saktësi në mantisë. Distanca relative ndërmjet numrave pranë njësisë është afërsisht $2^{-52}$. Madhësia
\begin{equation}
 u=\tfrac12\,\beta^{1-p}
\end{equation}
quhet njësia e rrumbullakimit. Modeli bazë është
\begin{equation}
 \operatorname{fl}(x\circ y)=(x\circ y)(1+\delta),\qquad |\delta|\le u,
\end{equation}
për një veprim $\circ$, me kusht që të mos ndodhë tejmbushje ose nënmbushje.

Numrat \texttt{Inf}, \texttt{-Inf} dhe \texttt{NaN} sinjalizojnë rezultate jashtë bashkësisë së numrave të fundmë ose veprime të papërcaktuara. Krahasimi i numrave realë nuk duhet të mbështetet verbërisht në barazinë bit për bit; përdoret një tolerancë që merr parasysh shkallën e problemit.

\section{Gabimi absolut dhe relativ}
Nëse $x$ është vlera e saktë dhe $\tilde x$ përafrimi,
\begin{equation}
 e_a=|\tilde x-x|,\qquad e_r=\frac{|\tilde x-x|}{|x|}.
\end{equation}
Gabimi relativ nuk ka kuptim kur $x=0$ dhe bëhet i paqëndrueshëm pranë zeros. Në vektorë përdoret një normë, për shembull
\begin{equation}
 e_r=\frac{\|\tilde{\vect x}-\vect x\|_2}{\|\vect x\|_2}.
\end{equation}

\section{Anulimi dhe humbja e shifrave}
Zbritja e dy numrave pothuaj të barabartë heq shifrat kryesore të përbashkëta dhe zmadhon gabimin relativ të mbetjes. Shprehja $\sqrt{x+1}-\sqrt{x}$ për $x\gg1$ llogaritet më mirë si
\begin{equation}
 \frac{1}{\sqrt{x+1}+\sqrt{x}}.
\end{equation}
Të dyja janë algjebrikisht të njëjta, por jo numerikisht të barasvlershme.

\begin{figure}[ht]
\centering\includegraphics[width=.75\textwidth]{02_anulimi_katastrofik.pdf}
\caption{Gabimi i shkaktuar nga anulimi kur $x$ bëhet i krahasueshëm me saktësinë e makinës.}
\end{figure}

\begin{kujdes}
Renditja e veprimeve ka rëndësi. Në një shumë me terma me madhësi shumë të ndryshme, termat e vegjël mund të humbasin. Mbledhja nga moduli më i vogël te më i madhi ose algoritmi Kahan redukton humbjen.
\end{kujdes}

\section{Kushtëzimi i problemit}
Kushtëzimi është veti e problemit matematik. Për $y=f(x)$, numri relativ i kushtëzimit lokal është
\begin{equation}
 \kappa_f(x)=\left|\frac{x f'(x)}{f(x)}\right|.
\end{equation}
Në sistemin $A\vect x=\vect b$,
\begin{equation}
 \frac{\|\delta\vect x\|}{\|\vect x\|}\lesssim \kappa(A)\frac{\|\delta\vect b\|}{\|\vect b\|},
 \qquad \kappa(A)=\|A\|\,\|A^{-1}\|.
\end{equation}
Një matricë me $\kappa\gg1$ është e keqkushtëzuar: perturbime të vogla në të dhëna mund të prodhojnë ndryshime të mëdha në zgjidhje.

\begin{figure}[ht]
\centering\includegraphics[width=.75\textwidth]{04_kushtzimi_matrices.pdf}
\caption{Rritja e numrit të kushtëzimit kur dy rreshta të matricës bëhen pothuaj linearë të varur.}
\end{figure}

\section{Gabimi i përparmë dhe i prapmë}
Gabimi i përparmë mat distancën ndërmjet rezultatit numerik dhe zgjidhjes së saktë. Analiza e prapme kërkon problemin pak të perturbur që zgjidhet saktësisht nga rezultati i llogaritur. Një algoritëm është i qëndrueshëm në kuptimin e prapëm nëse perturbimi i kërkuar është i vogël. Marrëdhënia konceptuale është
\begin{equation}
 \text{gabimi i përparmë}\ \lesssim\ \text{kushtëzimi}\times\text{gabimi i prapëm}.
\end{equation}
Prandaj, edhe algoritmi më i mirë nuk mund ta shpëtojë një problem intrinsikisht të keqkushtëzuar.

\section{Gabimet matricore}
Shumëzimi matricor kryen shumë mbledhje dhe shumëzime. Kufiri klasik për një produkt skalar me $n$ terma ka faktor
\begin{equation}
 \gamma_n=\frac{nu}{1-nu}\approx nu.
\end{equation}
Invertimi eksplicit i matricës për të zgjidhur $A\vect x=\vect b$ është zakonisht më i shtrenjtë dhe më pak i qëndrueshëm se faktorizimi i drejtpërdrejtë. Në kod përdoret operatori i zgjidhjes, jo $A^{-1}\vect b$.

\begin{shembull}[Shuma e energjive]
Nëse energjia totale merret si diferencë e dy energjive të mëdha, gabimi relativ mund të jetë i lartë edhe kur secila energji është llogaritur me shumë shifra. Një formulim që llogarit drejtpërdrejt madhësinë e vogël mund të jetë numerikisht më i sigurt.
\end{shembull}

\begin{lidhje}
Në laborator maten $u$, anulimi dhe ndikimi i rendit të mbledhjes. Pastaj ndërtohet një familje sistemesh $2\times2$ dhe krahasohet gabimi i zgjidhjes me $\kappa(A)$.
\end{lidhje}

\section*{Pyetje kontrolli}
\begin{enumerate}
\item Dalloni një problem të keqkushtëzuar nga një algoritëm i paqëndrueshëm.
\item Pse $A^{-1}\vect b$ nuk është mënyra e rekomanduar për të zgjidhur sistemin?
\item Jepni një reformulim që shmang anulimin katastrofik.
\end{enumerate}
\section{Algoritmet për diagnostikimin e gabimeve numerike}
Qëllimi praktik është të dallohen katër burime: gabimi i të dhënave, kushtëzimi i problemit, gabimi i algoritmit dhe rrumbullakimi. Një diagnozë e mirë ndryshon vetëm një faktor në çdo provë.

\subsection{Matja e saktësisë së makinës}
\begin{lstlisting}[language=Python,caption={Vlerësimi empirik i distancës pranë njësisë.}]
eps = 1.0
while 1.0 + eps/2.0 != 1.0:
    eps /= 2.0
print(eps, np.finfo(float).eps)
\end{lstlisting}
Vlera e marrë është epsilon-i i makinës sipas konventës së NumPy. Njësia e rrumbullakimit është afërsisht gjysma e saj.

\subsection{Shuma Kahan}
Në mbledhjen e shumë termave, një kompensim ruan pjesën e humbur nga rrumbullakimi.
\begin{lstlisting}[language=Python,caption={Mbledhja e kompensuar Kahan.}]
def shuma_kahan(x):
    s = 0.0
    c = 0.0
    for value in x:
        y = value - c
        t = s + y
        c = (t - s) - y
        s = t
    return s

x = np.r_[1e16, np.ones(1_000_000), -1e16]
print(np.sum(x), shuma_kahan(x))
\end{lstlisting}
Ky test duhet përsëritur me renditje të ndryshme. Rezultati tregon se asociativiteti algjebrik nuk ruhet plotësisht në pikë të lëvizshme.

\subsection{Kushtëzimi dhe mbetja}
\begin{lstlisting}[language=Python,caption={Krahasimi i gabimit me mbetjen.}]
def diagnostiko(A, b, x_ref):
    x = np.linalg.solve(A, b)
    mbetja = np.linalg.norm(b - A @ x) / np.linalg.norm(b)
    gabimi = np.linalg.norm(x - x_ref) / np.linalg.norm(x_ref)
    return np.linalg.cond(A), mbetja, gabimi

for e in [1e-2, 1e-6, 1e-10]:
    A = np.array([[1.0, 1.0], [1.0, 1.0 + e]])
    x_ref = np.array([1.0, 1.0])
    b = A @ x_ref
    print(e, diagnostiko(A, b, x_ref))
\end{lstlisting}
Mbetja mund të mbetet e vogël edhe kur gabimi i zgjidhjes rritet. Prandaj raportohen të dyja bashkë me $\kappa(A)$.

\subsection{Algoritmi i auditimit}
\begin{enumerate}
\item Përcaktohet një zgjidhje reference me saktësi të lartë ose analitike.
\item Ndryshohet hapi, toleranca ose rendi i veprimeve.
\item Vizatohet gabimi në shkallë logaritmike.
\item Kontrollohet pjerrësia e pritur e konvergjencës.
\item Përsëritet me të dhëna pak të perturbura për të matur kushtëzimin.
\end{enumerate}

Terminologjia \emph{gabim i përparmë}, \emph{gabim i prapëm}, \emph{kushtëzim} dhe \emph{qëndrueshmëri} përdoret sipas traditës së analizës numerike në shqip, veçanërisht në tekstet e Hanellit--Osmanit dhe Hoxhës.

Për trajtimin në shqip të gabimeve, sistemeve lineare dhe qëndrueshmërisë janë përdorur \cite{hanelli,hoxha,datja}; analiza moderne e qëndrueshmërisë krahasohet me \cite{higham}.

\subsection{Modeli standard i rrumbullakimit dhe përhapja e gabimit}
Le të jetë $\operatorname{fl}(x)$ numri i përfaqësueshëm më i afërt me $x$. Për numrat normalë dhe rrumbullakimin drejt më të afërtit shkruhet
\begin{equation}
 \operatorname{fl}(x)=x(1+\delta),\qquad |\delta|\le u.
\end{equation}
Ky model është relativ dhe prandaj nuk përshkruan drejtpërdrejt sjelljen pranë nënmbushjes. Për një varg veprimesh, çdo hap fut një faktor të ri $(1+\delta_i)$. Prodhimi plotëson
\begin{equation}
 \prod_{i=1}^{n}(1+\delta_i)=1+\theta_n,qquad
 |\theta_n|\le\gamma_n=\frac{nu}{1-nu},
\end{equation}
kur $nu<1$. Formula $\gamma_n\simeq nu$ është e dobishme, por nuk duhet interpretuar si parashikim i saktë i gabimit: ajo është kufi në rastin më të pafavorshëm.

Për $f(x_1,\ldots,x_p)$, linearizimi jep
\begin{equation}
 \Delta f\simeq\sum_{j=1}^{p}\frac{\partial f}{\partial x_j}\Delta x_j.
\end{equation}
Nëse perturbimet trajtohen si të pavarura me devijime standarde $\sigma_j$, përhapja probabilitare e rendit të parë është
\begin{equation}
 \sigma_f^2\simeq\sum_{j=1}^{p}\left(\frac{\partial f}{\partial x_j}\right)^2\sigma_j^2,
\end{equation}
ndërsa për hyrje të korreluara shfaqen termat $2(\partial_i f)(\partial_j f)\operatorname{Cov}(x_i,x_j)$. Analiza deterministe e rrumbullakimit dhe përhapja probabilitare e pacaktueshmërisë nuk janë e njëjta gjë, megjithëse përdorin të njëjtin mjet të linearizimit.

\subsection{Kushtëzimi absolut dhe relativ}
Për një hartë $f:\R^n\to\R^m$, perturbimi i rendit të parë është $\Delta\vect y\simeq J_f(\vect x)\Delta\vect x$. Numri absolut i kushtëzimit në normat e zgjedhura është
\begin{equation}
 \kappa_{\mathrm{abs}}(\vect x)=\|J_f(\vect x)\|,
\end{equation}
kurse numri relativ mund të shkruhet
\begin{equation}
 \kappa_{\mathrm{rel}}(\vect x)=
 \frac{\|J_f(\vect x)\|\,\|\vect x\|}{\|f(\vect x)\|}.
\end{equation}
Kjo formulë tregon pse kushtëzimi rritet pranë një rezultati zero: një ndryshim absolut i moderuar bëhet ndryshim relativ i madh. Për sistemin linear, nga
\begin{equation}
 (A+\Delta A)(\vect x+\Delta\vect x)=\vect b+\Delta\vect b
\end{equation}
dhe neglizhimi i produkteve të rendit të dytë merret
\begin{equation}
 A\Delta\vect x\simeq\Delta\vect b-\Delta A\vect x.
\end{equation}
Prandaj
\begin{equation}
 \frac{\|\Delta\vect x\|}{\|\vect x\|}
 \lesssim\kappa(A)\left(
 \frac{\|\Delta A\|}{\|A\|}+\frac{\|\Delta\vect b\|}{\|\vect b\|}
 \right),
\end{equation}
me kusht që perturbimi i matricës të jetë mjaft i vogël.

\subsection{Mbetja dhe kufiri i gabimit të zgjidhjes}
Për një përafrim $\widehat{\vect x}$, mbetja është $\vect r=\vect b-A\widehat{\vect x}$. Gabimi i vërtetë $\vect e=\vect x-\widehat{\vect x}$ plotëson $A\vect e=\vect r$, kështu që
\begin{equation}
 \|\vect e\|\le\|A^{-1}\|\,\|\vect r\|.
\end{equation}
Duke normalizuar dhe përdorur $\|\vect b\|\le\|A\|\|\vect x\|$ merret kufiri praktik
\begin{equation}
 \frac{\|\vect e\|}{\|\vect x\|}
 \le \kappa(A)\frac{\|\vect r\|}{\|\vect b\|}.
\end{equation}
Pra mbetja e vogël është e nevojshme, por jo e mjaftueshme për saktësi të lartë kur $\kappa(A)$ është e madhe.

\begin{lstlisting}[language=Python,caption={Raport diagnostik për një sistem linear.}]
import numpy as np

def raport_sistemi(A, b, x_hat):
    r = b - A @ x_hat
    mbetje_relative = np.linalg.norm(r) / np.linalg.norm(b)
    kappa = np.linalg.cond(A)
    kufi_gabimi = kappa * mbetje_relative
    return {"kappa": kappa,
            "mbetje_relative": mbetje_relative,
            "kufi_i_perafert": kufi_gabimi}
\end{lstlisting}
Funksioni nuk pretendon se kufiri është gabimi real; ai jep një alarm të arsyetuar. Në MATLAB/Octave përdoren \texttt{norm}, \texttt{cond} dhe prodhimi \texttt{A*x\_hat}.

\subsection{Përmbledhje dhe ushtrime}
Gabimi i të dhënave, kushtëzimi i problemit dhe qëndrueshmëria e algoritmit duhet të analizohen veçmas. Një raport numerik pa mbetje, kushtëzim dhe studim të saktësisë është i paplotë.
\begin{enumerate}
\item Nxirrni numrin relativ të kushtëzimit të $f(x)=1/(1-x)$ dhe interpretojeni pranë $x=1$.
\item Provoni kufirin $|\theta_n|\le\gamma_n$ për $n=2$ dhe diskutoni rolin e termit $u^2$.
\item Krahasoni shumën naive, shumën nga termi më i vogël dhe shumën Kahan për një varg me shkallë të ndryshme.
\item Ndërtoni matrica me numër kushtëzimi të kontrolluar nga vlerat singulare dhe studioni gabimin përpara.
\end{enumerate}

\subsection{Shuma, produkti skalar dhe gabimet e grumbulluara}
Shuma $s_n=\sum_{i=1}^{n}x_i$ duket operacion elementar, por është një nga burimet më të zakonshme të humbjes së saktësisë. Në mbledhjen naive,
\begin{equation}
 \widehat s_k=\operatorname{fl}(\widehat s_{k-1}+x_k),
\end{equation}
secili hap rrumbullakon rezultatin e pjesshëm. Analiza standarde jep
\begin{equation}
 |\widehat s_n-s_n|\le\gamma_{n-1}\sum_{i=1}^{n}|x_i|,
\end{equation}
ku $\gamma_{n-1}\simeq(n-1)u$. Kufiri varet nga shuma e moduleve, jo nga moduli i shumës. Kur ka anulim të fortë, $|s_n|\ll\sum|x_i|$, gabimi relativ mund të jetë i madh edhe pse gabimi absolut respekton kufirin.

Mbledhja në çift formon një pemë binare dhe ka thellësi $\lceil\log_2n\rceil$ në vend të $n-1$; kufiri përkatës rritet afërsisht me $u\log_2n$. Kjo është e rëndësishme në llogaritjen paralele, ku rendi i reduktimit ndryshon rezultatin bit për bit. Algoritmi Kahan ruan një kompensim për pjesën e humbur dhe shpesh arrin gabim të krahasueshëm me një rrumbullakim të vetëm, megjithëse nuk është ilaç universal për të gjitha vargjet.

Produkti skalar $\vect x^T\vect y$ kombinon shumëzime dhe mbledhje. Kufiri tipik është
\begin{equation}
 |\operatorname{fl}(\vect x^T\vect y)-\vect x^T\vect y|
 \le\gamma_n |\vect x|^T|\vect y|.
\end{equation}
Ky rezultat përhapet te shumëzimi matricor element për element. Në simulime të mëdha, ndryshimi i bibliotekës BLAS, numrit të threads ose rendit të ndarjes mund të ndryshojë bitat e fundit pa ndryshuar vlefshmërinë shkencore. Riprodhueshmëria duhet të përcaktojë nëse kërkohet identitet bit për bit apo vetëm pajtim brenda një tolerance.

\subsubsection{Kushtëzimi komponentor}
Numri klasik $\kappa(A)=\|A\|\|A^{-1}\|$ mat perturbime normore. Në disa probleme, hyrjet kanë pacaktueshmëri relative komponent më komponent. Atëherë një matricë me elemente në shkallë shumë të ndryshme mund të duket keq në një normë, edhe pse struktura fizike e perturbimeve është më e kufizuar. Analiza komponentore përdor $|\Delta A|\le\varepsilon|A|$ dhe $|\Delta b|\le\varepsilon|b|$. Ky dallim është me rëndësi në rrjete rezistive, sisteme bilanci dhe matrica që përmbajnë koeficientë me njësi të ndryshme.

Shkallëzimi me matrica diagonale $D_rAD_c$ mund të barazojë madhësitë e rreshtave dhe kolonave. Ai nuk ndryshon zgjidhjen fizike pasi ndryshoret rikthehen, por mund të përmirësojë sjelljen e eliminimit dhe kuptimin e tolerancave. Shkallëzimi nuk e shëron një varësi reale lineare; ai vetëm shmang që njësitë e papërshtatshme të duken si keqkushtëzim.

\subsubsection{Aritmetika e intervaleve si kontroll}
Në aritmetikën e intervaleve, një madhësi ruhet si $[a,b]$ që garanton përfshirjen e vlerës së saktë. Veprimet rrumbullakohen nga jashtë. Për shembull,
\begin{equation}
 [a,b]+[c,d]=[a+c,b+d].
\end{equation}
Avantazhi është një kufi i garantuar; kufizimi është mbivlerësimi nga varësia. Nëse $x\in[0,1]$, shprehja $x-x$ vlerësohet si $[-1,1]$, megjithëse është saktësisht zero, sepse dy paraqitjet e $x$ trajtohen si të pavarura. Aritmetika e intervaleve është e dobishme për verifikime lokale dhe kufij, por nuk zëvendëson analizën e kushtëzimit ose modelin probabilitar të pacaktueshmërisë.

